\subsection{A primitive rank-zero cover removes the elliptic exception}
\label{sec:geometric-primitive-all-genera}

\paragraph{Finding A18 (geometric primitive criterion in all genera).}
Let $C/k$ be a smooth projective geometrically connected curve of
genus $g\ge2$ over a finite field of characteristic two.
Suppose that, over a finite extension of constants, there is a
geometrically connected cyclic \'etale cover $D\to C$ of odd degree
$n>1$ whose primitive conductor-$n$ Jacobian factor has geometric
$2$-rank zero.  Then the finite-field point-section separation
conclusion in the covering formulation of A1 holds for $C$.
There is no hypothesis excluding a separable double map from $C$
to a genus-one curve, and no restriction that $n$ be a prime power.

More precisely, with constant deck group $G=\langle\sigma\rangle$,
write
\[
 A=\prod_{\ell\mid n}(1-\sigma^{n/\ell}),\qquad
 P_n=\operatorname{im}(A:J_D\longrightarrow J_D),
 \tag{GP.1}
\]
where the product runs over the distinct prime divisors of $n$.
Under a putative collision, if $P_n(L)[2^\infty]=0$ on a pointed
model over a finite field $L$, the exceptional outcome of the
projector argument necessarily satisfies
\[
                         f(P_n)=r\varphi(n)>0
 \tag{GP.2}
\]
for an integer $r\ge1$.  Here $r$ is the number of geometric
ramification points of a resulting separable degree-two map
$C_{\overline L}\to E$ with $g(E)=1$.
For genus two, $r=1$, and the entire nontrivial new factor of
$D/C$ in this exceptional case is ordinary.
The geometric rank-zero hypothesis therefore eliminates this
exception.  A rational torsion condition alone is not being
replaced by a geometric one.

\begin{proof}
\emph{1. The projector reduction.}
Extend the constants so that the proposed distinct colliding
points $P,Q$, the cover, its cyclic deck group and a chosen lift
$P'\in D(L)$ are all defined over one finite field $L$.
Geometric $2$-rank zero implies
$P_n(L)[2^\infty]=0$ over every such extension.
The covering formulation of a collision gives a lift $Q'$ over
$Q$ with
\[
                       \delta=[Q'-P']\in J_D(L)[2^\infty].
\]
Hence $A\delta=0$.

We recall the complete geometric reduction from A15,
Steps~2--4; those steps do not use the rank or genus-two
hypothesis on the base except for the lower bound on $g(D)$.
Put $s=\omega(n)$ and $d=2^s$.
The $2^s$ subset products of the elements $\sigma^{n/\ell}$
are distinct, because they lie in independent prime-primary
factors of the cyclic group.  The two deck orbits of $P',Q'$
are free and disjoint.  Thus the expansion of the principal
divisor $A(Q'-P')$ has precisely $d$ distinct simple zeros and
$d$ distinct simple poles.  The resulting function $f$ has
separable degree $d$.

Over $\overline L$, put
\[
 K=\overline L(D),\qquad
 M=\overline L(\gamma f:\gamma\in G),\qquad e=[K:M].
\]
The Castelnuovo--Severi induction in A15 gives
\[
                 e\mid d,\qquad g(M)\le(d/e-1)^2.
 \tag{GP.3}
\]
The exact published input is Poonen,
\emph{Gonality of modular curves in characteristic $p$},
Proposition~2.1, Lemma~2.2 and Corollary~2.3,
author p.~3 \cite{Poonen2007Gonality}.  The induction adjoins functions of degree $d$;
if the previous index is $d/u$ and the adjoining degree is $a$,
the new index is $d/(au)$ and its genus is at most
$a(u-1)^2+(a-1)(au-1)\le(au-1)^2$.
For every odd $n>1$ one has $n>(d-1)^2$:
for $s=1,2$ use $n\ge3,15$, and for $s\ge3$ use
$n\ge3\cdot5^{s-1}>4^s$.
Since $g(D)=n(g-1)+1\ge n+1$, (GP.3) excludes $e=1$.
Thus $e>1$ is a power of two.

Let $Z$ be the smooth projective curve with function field $M$.
The induced $G$-action on $Z$ is faithful: a kernel would have
order dividing both $e$ and $n$.
It is free: a stabilizer containing an element of odd prime
order $\ell$ would partition a fibre divisor of $D\to Z$
of degree $e$ into orbits of length $\ell$, with constant
multiplicities on each orbit, forcing $\ell\mid e$.
Consequently
\[
                 g(Z)-1=n\bigl(g(Z/G)-1\bigr).
\]
Together with $g(Z)\le(d/e-1)^2<n$, this forces
$g(Z)=g(Z/G)=1$.

Write $f=h\circ v$ for $v:D\to Z$.
A fibre above a zero of $h$ is reduced, since the zero divisor
of $f$ is reduced.  It has $e$ points, all in the two original
deck orbits.  Each of these orbits maps injectively to $Z$,
because the action on $Z$ is free.  Hence $e\le2$, so $e=2$.
The extension $K/M$ is separable and quadratic, with involution
$\tau$.  It is separable because it is an intermediate extension
of the separable extension $K/\overline L(f)$.
Since $G$ preserves $M$, it normalizes $\langle\tau\rangle$;
the unique nontrivial element is fixed by conjugation, so
$\tau$ commutes with $G$.

\emph{2. The exceptional diagram is an actual pullback.}
Set
\[
 F=K^G=\overline L(C),\qquad
 E_0=M^G=K^{G\times\langle\tau\rangle}.
\]
Let $E$ be the smooth projective curve with function field $E_0$.
We already know that $g(E)=1$.
The group $G\times\langle\tau\rangle$ acts faithfully on $K$,
and has order $2n$, so
$K/E_0$ is Galois of that degree.  The fields fixed by
$G$ and by $\langle\tau\rangle$ have intersection $E_0$ and
compositum $K$.  Equivalently,
\[
       [F:E_0]=2,\qquad [M:E_0]=n,\qquad FM=K.
 \tag{GP.4}
\]
In particular $C\to E$ is separable of degree two and
$Z\to E$ is the cyclic \'etale cover of degree $n$ already
obtained from the free $G$-action.
The fibre product $C\times_E Z$ is connected by (GP.4).
It is smooth, since its projection to $C$ is the base change
of an \'etale morphism.  Its function field is $K$, so it is
the curve $D$ itself.

For each divisor $m\mid n$, let $H_m$ be the subgroup of $G$
of order $n/m$, and set
\[
       Z_m=Z/H_m,\qquad D_m=D/H_m.
\]
The same field calculation gives
\[
        D_m=C\times_E Z_m,\qquad
        \deg(Z_m/E)=m,\qquad \deg(D_m/C)=m.
 \tag{GP.5}
\]
Every $Z_m\to E$ is \'etale between genus-one curves.
After choosing compatible origins over $\overline L$ it is
an isogeny of odd degree.  Such an isogeny induces an
isomorphism on the geometric $2$-primary Tate modules:
its kernel has odd order, and multiplication by its degree
is invertible on every $2$-primary group.  Therefore
\[
                              f(Z_m)=f(E).
 \tag{GP.6}
\]

\emph{3. Ramification and the exact ranks of all intermediate covers.}
Let $r$ be the number of geometric ramification points of
$C\to E$.  Riemann--Hurwitz gives different of degree
$2g-2>0$, so $r\ge1$.  This is a quadratic Galois extension
in characteristic two.  Its base change by the \'etale map
$Z_m\to E$ has exactly $mr$ ramification points: above each
branch point are $m$ distinct points of $Z_m$, and \'etale
base change preserves the local ramification data.
The Deuring--Shafarevich formula for this degree-two cover,
together with (GP.6), gives
\[
             f(D_m)-1=2\bigl(f(Z_m)-1\bigr)+mr,
 \qquad
             f(D_m)=2f(E)-1+mr.
 \tag{GP.7}
\]
The formula counts the geometric ramification points, since
each has ramification index two.  It does not replace that
count by the degree of the different.
For the exact primary locator of the formula, see
Anbar--Beelen--Nguyen,
\emph{A new tower with good $p$-rank meeting Zink's bound},
Theorem~1, equation~(1), author p.~3 \cite{AnbarBeelenNguyenP}.
Its hypotheses are a Galois $p$-group extension of function
fields over an algebraically closed field, as here.

\emph{4. Conductor decomposition forces positive primitive rank.}
Consider
$V=T_2(J_D)\otimes_{\mathbf Z_2}\mathbf Q_2$,
whose dimension is the geometric $2$-rank of $D$.
The odd group $G$ acts semisimply on $V$.
Pullback and norm for $D\to D_m$ have composition
$[n/m]$ in one direction and the sum of the subgroup
$H_m$ in the other.  Since $n/m$ is odd, they identify
$T_2(J_{D_m})\otimes\mathbf Q_2$ with $V^{H_m}$.

After extending the coefficient field to split $G$, decompose
$V$ according to the exact orders of its characters.
Let $w_j$ be the total dimension of the characters of exact
order $j$, for $j\mid n$.  The invariants just identified give
\[
                           f(D_m)=\sum_{j\mid m}w_j.
 \tag{GP.8}
\]
On character $i$ modulo $n$, the endomorphism $A$ in (GP.1)
acts by
$\prod_{\ell\mid n}(1-\zeta_\ell^i)$.
This is nonzero exactly when $\gcd(i,n)=1$.
Thus its image has the primitive character part as its
$2$-adic Tate module up to isogeny, and
$f(P_n)=w_n$.
M\"obius inversion of (GP.8), using (GP.7) and $n>1$, yields
\[
\begin{aligned}
 f(P_n)
   &=\sum_{m\mid n}\mu(n/m)\bigl(2f(E)-1+mr\bigr)\\
   &=r\sum_{m\mid n}\mu(n/m)m
     =r\varphi(n)>0.
\end{aligned}
 \tag{GP.9}
\]
This contradicts the geometric rank-zero hypothesis.
Hence the proposed collision cannot occur.
\end{proof}

\paragraph{Genus-two consequence and the rational distinction.}
For $g=2$, the different of a separable quadratic map to a
genus-one curve has degree two.  Each ramified point in
characteristic two has different exponent at least two.
Thus $r=1$, with a single ramified point of different exponent
two.  Formula (GP.7) gives
$f(D)-f(C)=n-1$, equal to $g(D)-g(C)$.
The pullback/norm decomposition of $J_D$ into $J_C$ and its
new factor, up to an odd-degree isogeny, therefore shows that
the whole new factor is ordinary.  In particular its primitive
factor is ordinary of dimension $\varphi(n)$.

Consequently, even the rational condition
$P_n(L)[2^\infty]=0$ suffices on a genus-two base if one also
knows that $P_n$ is \emph{not ordinary}: a collision would force
it to be ordinary.  Without this extra geometric condition,
the rational-only statement of A15 retains its stated
elliptic-quotient qualification.  In all genera, (GP.9)
describes the positive primitive rank forced by the
exceptional case.

The theta calculation of A15 now gives a smaller geometric
gate on \emph{every} genus-two curve in characteristic two,
of any $2$-rank: an exact odd-order theta point with
$\langle2,-1\rangle=(\mathbf Z/n)^\times$ is sufficient,
and so is the exact same-subgroup orbit-coverage condition
there.  The arithmetic family $n=3q$ of A17 is therefore
available as a family of candidate orders on ordinary
bielliptic bases as well.
This does not assert incidence for every curve.  The
one-dimensional eigenspace argument establishing that
theta criterion remains specific to genus two; A18's
all-genus assertion assumes a primitive factor of geometric
rank zero directly.
